\section{Clauses, correlation, and restrictions}\label{sec:preliminaries}

We record the clause and restriction conventions needed for both structural
lemmas.

For a finite coordinate set \(V\), the cube \(\bits^V\) carries the uniform
probability measure.  Expectations without a specified measure use this
measure.  A Boolean predicate is identified with its \(0/1\) indicator.
Its associated sign is \(+1\) on accepted inputs and \(-1\) otherwise.

A \emph{clause} is an OR of literals, and a \emph{CNF} is an AND of
clauses.  Tautological and constant-true clauses can be removed.
Each remaining nonconstant clause uses distinct variables and contains
no constants.  Its \emph{scope} is the set of these variables, and its
\emph{width} is the size of its scope.  It is violated on exactly one
assignment to its scope.  A constant-false clause has empty scope and
width zero.  The empty CNF is the constant-true function.

For an integer \(b\ge1\), let \(\mathcal C_b(V)\) be the set of CNF
indicators on \(\bits^V\) with clause widths at most \(b\).
There is no bound on the number of clauses.  For
\(G:\bits^V\to\R\), define
\begin{equation}\label{eq:correlation-definition}
 \corr_b(G)=\max_{J\in\mathcal C_b(V)}
               \abs{\E_{z\in\bits^V}G(z)J(z)}.
\end{equation}
The maximum exists: there are finitely many distinct clause indicators,
and repetitions do not change a CNF's function.  Both constant functions
belong to \(\mathcal C_b(V)\).

A \(p\)-random restriction of \([d]=\{1,\ldots,d\}\), for \(0<p<1\),
is sampled in two stages.  Include each coordinate independently in
a live set \(T\) with probability \(p\); then choose an independent
uniform assignment \(\rho\) to \([d]\setminus T\).
For \(G:\bits^d\to\R\), write \(G_{T,\rho}\) for the function remaining
on \(\bits^T\).  Filling the live coordinates with uniform bits produces
a uniform point of \(\bits^d\).  Consequently, for any real-valued
functions \(G,H\) on \(\bits^d\),
\begin{equation}\label{eq:restriction-expectation}
 \E_xG(x)H(x)
  =\E_{T,\rho}\E_zG_{T,\rho}(z)H_{T,\rho}(z).
\end{equation}
We will also use a \emph{cylinder}: prescribing a pattern on
\(U\subseteq T\) and leaving \(T\setminus U\) arbitrary.
Its indicator is a conjunction of unit clauses, and its measure on
the live cube is \(2^{-|U|}\).
